\chapter{DISCUSIÓN}

En este capítulo se analizan los resultados obtenidos en la experimentación con cuatro modelos de detección de crimen financiero, evaluados tanto en su desempeño original como después de ser sometidos a cuatro tipos de ataques adversariales sobre tres configuraciones de datasets. Los resultados muestran una vulnerabilidad generalizada de los modelos ante estos ataques, reflejada principalmente en la tasa de evasión y en la norma L\textsubscript{0}, que mide cuántas features fueron modificadas para lograr que una transacción fraudulenta sea clasificada como legítima.

% ─────────────────────────────────────────────────────────────
\section{Vulnerabilidad de los Modelos}
% ─────────────────────────────────────────────────────────────

Los resultados presentados en las Tablas~\ref{tab:cc_results}, \ref{tab:aml_nosmote} y \ref{tab:aml_smote} demuestran que la mayoría de los modelos evaluados sufren una degradación considerable de su capacidad de detección al ser sometidos a ataques adversariales. Sin embargo, el grado de vulnerabilidad varía significativamente según el modelo, el dataset y el tipo de ataque.

En el dataset de Credit Card, Regresión Logística y MLP son los más afectados: HopSkipJump y SquareAttack logran evasión superior al 99\,\% contra ambos modelos, dejándolos prácticamente inoperativos para detectar fraude. XGBoost también muestra alta vulnerabilidad frente a HopSkipJump y BoundaryAttack (evasión del 100\,\%). En contraste, LSTM-Att. presenta mayor resistencia frente a CaFA (23.7\,\%) y SquareAttack (22.6\,\%), aunque sigue siendo vulnerable a HopSkipJump (77.9\,\%) cuando el atacante está dispuesto a modificar prácticamente todas las features.

En AMLworld sin SMOTE, el desbalance extremo de clases (1:979) impide que MLP y LSTM-Att. aprendan los patrones de fraude en primer lugar, alcanzando FNR\,=\,100\,\% en baseline y quedando excluidos de la evaluación adversarial. Los dos modelos evaluables —Regresión Logística y XGBoost— muestran una vulnerabilidad crítica: CaFA, BoundaryAttack y SquareAttack logran el 100\,\% de evasión modificando menos del 15\,\% de las features. Esto indica que los modelos entrenados sin balanceo de clases aprenden una región de fraude tan pequeña y frágil que cualquier perturbación mínima es suficiente para sacar una muestra de esa región.

El resultado más destacado aparece en AMLworld con SMOTE: \textbf{LSTM-Att.(S) resiste prácticamente todos los ataques de caja negra} (HopSkipJump: 0.0\,\%, BoundaryAttack: 0.5\,\%, SquareAttack: 0.2\,\%). Durante la discusión con la asesora surgió la pregunta de por qué este modelo es el más robusto. La respuesta tiene dos partes. La primera es arquitectónica: la LSTM procesa las features en secuencia a través de celdas con compuertas y el mecanismo de atención aprende cuáles features importan más, generando una frontera de decisión irregular que los ataques de caja negra no logran aproximar con sus consultas. La segunda razón es más sutil: LSTM-Att.(S) tiene un Recall del 96.04\,\% pero un F1 del 0.25\,\%, lo que indica que el modelo está fuertemente sesgado a clasificar todo como fraude. Un modelo así de agresivo en la detección es paradójicamente muy difícil de engañar —el atacante necesita convencerlo de que una transacción fraudulenta es legítima, pero el modelo ya está predispuesto a clasificar casi todo como fraude. Esta observación revela una tensión importante: \textbf{robustez adversarial y utilidad operacional no siempre van de la mano}. El modelo más resistente a los ataques resulta ser el menos útil en producción por su alta tasa de falsas alarmas.

En el extremo opuesto, MLP\,(S) es el modelo más vulnerable en AMLworld con SMOTE: los cuatro ataques logran evasión superior al 95\,\% modificando menos del 24\,\% de las features.

% ─────────────────────────────────────────────────────────────
\section{Impacto de los Ataques Adversariales}
% ─────────────────────────────────────────────────────────────

El análisis del porcentaje de features modificadas (L\textsubscript{0}\,/\,n) y de la Zona de Peligro proporciona una visión más detallada del impacto real de cada ataque. Se observan diferencias notables entre los cuatro ataques evaluados tanto en efectividad como en eficiencia.

CaFA, al ser un ataque de caja blanca con acceso al gradiente del modelo, logra los menores costos estructurales: en promedio modifica solo el 16.4\,\% de las features para lograr evasión, siendo el ataque más eficiente en términos de perturbación mínima. Esto lo posiciona como la amenaza más preocupante en escenarios donde el atacante tiene conocimiento del modelo —por ejemplo, en ataques internos o tras una filtración del sistema. BoundaryAttack, en cambio, presenta la mayor evasión media entre los ataques de caja negra (80.7\,\%), siendo el más consistentemente efectivo a lo largo de los tres datasets.

HopSkipJump muestra el comportamiento más variable. En Credit Card es muy efectivo (evasión cercana al 100\,\% en tres de los cuatro modelos), pero en AMLworld sin SMOTE falla casi completamente (5.2\,\% y 2.2\,\% de evasión). La razón es que HopSkipJump estima el gradiente de la frontera de decisión a partir de consultas binarias al modelo, asumiendo que esa frontera es aproximadamente regular. En AMLworld sin SMOTE, el desbalance extremo genera una frontera muy asimétrica que el algoritmo no logra aproximar, y sus consultas no convergen a una dirección de ataque útil. Los ataques que no dependen de esta estimación —CaFA tiene el gradiente real, BoundaryAttack hace pasos ortogonales, SquareAttack perturba subconjuntos aleatorios— no tienen este problema.

Respecto al efecto del SMOTE sobre la vulnerabilidad, los resultados muestran que el rebalanceo es condición necesaria para que MLP y LSTM-Att. puedan ser evaluados adversarialmente, pero también amplía la superficie de ataque: con SMOTE, el 73\,\% de las combinaciones caen en Zona de Peligro y el 67\,\% en Zona Crítica, frente al 60\,\% y 20\,\% respectivamente en Credit Card. SMOTE genera muestras sintéticas que densifican la región de fraude, lo cual permite a los modelos aprender, pero al mismo tiempo crea más rutas por las que una perturbación puede sacar una muestra de esa región con pocos cambios.

Finalmente, la comparación entre datasets revela que AMLworld es estructuralmente más vulnerable que Credit Card. Las features de AMLworld son más interpretables —montos, tipos de transacción, entidades— y tienen efectos más directos sobre la predicción, facilitando que los ataques encuentren perturbaciones efectivas. Credit Card, con sus 29 features derivadas de PCA, ofrece un espacio de búsqueda más complejo donde las perturbaciones tienen efectos menos predecibles sobre la decisión del modelo.
